\documentclass[10pt,a4paper]{amsart}
\usepackage{amsmath,amssymb,amsthm,graphicx}
% Retained minimal preamble for the Stacks Project chapter more-algebra.tex (More on Algebra).
% The chapter file pinned-original-source.tex carries only \input{preamble} on its first line, so no
% preamble text lives inside the chapter itself.  The Stacks Project uses one single shared preamble for
% every chapter; the snapshot stages it identically as sources/stacks-chow/preamble.tex and
% sources/stacks-etale-cohomology/preamble.tex.  The block below is copied mechanically, line for line,
% from sources/stacks-chow/preamble.tex lines 14-18 (the Xy-pic setup) and 154-217 (the native theorem
% declarations together with the mathematical macros).  Nothing was added, removed or redefined; the
% wrapper class and the amsmath/amssymb/amsthm/graphicx base packages are supplied by the assembler.
% Chapter: more-algebra.tex -- More on Algebra.
% For commutative diagrams we use Xy-pic
\usepackage[all]{xy}

% We use 2cell for 2-commutative diagrams.
\xyoption{2cell}

% Theorem environments.
%
\theoremstyle{plain}
\newtheorem{theorem}[subsection]{Theorem}
\newtheorem{proposition}[subsection]{Proposition}
\newtheorem{lemma}[subsection]{Lemma}

\theoremstyle{definition}
\newtheorem{definition}[subsection]{Definition}
\newtheorem{example}[subsection]{Example}
\newtheorem{exercise}[subsection]{Exercise}
\newtheorem{situation}[subsection]{Situation}

\theoremstyle{remark}
\newtheorem{remark}[subsection]{Remark}
\newtheorem{remarks}[subsection]{Remarks}

\numberwithin{equation}{subsection}

% Macros
%
\def\lim{\mathop{\mathrm{lim}}\nolimits}
\def\colim{\mathop{\mathrm{colim}}\nolimits}
\def\Spec{\mathop{\mathrm{Spec}}}
\def\Hom{\mathop{\mathrm{Hom}}\nolimits}
\def\Ext{\mathop{\mathrm{Ext}}\nolimits}
\def\SheafHom{\mathop{\mathcal{H}\!\mathit{om}}\nolimits}
\def\SheafExt{\mathop{\mathcal{E}\!\mathit{xt}}\nolimits}
\def\Sch{\mathit{Sch}}
\def\Mor{\mathop{\mathrm{Mor}}\nolimits}
\def\Ob{\mathop{\mathrm{Ob}}\nolimits}
\def\Sh{\mathop{\mathit{Sh}}\nolimits}
\def\NL{\mathop{N\!L}\nolimits}
\def\CH{\mathop{\mathrm{CH}}\nolimits}
\def\proetale{{pro\text{-}\acute{e}tale}}
\def\etale{{\acute{e}tale}}
\def\QCoh{\mathit{QCoh}}
\def\Ker{\mathop{\mathrm{Ker}}}
\def\Im{\mathop{\mathrm{Im}}}
\def\Coker{\mathop{\mathrm{Coker}}}
\def\Coim{\mathop{\mathrm{Coim}}}

% Boxtimes
%
\DeclareMathSymbol{\boxtimes}{\mathbin}{AMSa}{"02}

%
% Macros for moduli stacks/spaces
%
\def\QCohstack{\mathcal{QC}\!\mathit{oh}}
\def\Cohstack{\mathcal{C}\!\mathit{oh}}
\def\Spacesstack{\mathcal{S}\!\mathit{paces}}
\def\Quotfunctor{\mathrm{Quot}}
\def\Hilbfunctor{\mathrm{Hilb}}
\def\Curvesstack{\mathcal{C}\!\mathit{urves}}
\def\Polarizedstack{\mathcal{P}\!\mathit{olarized}}
\def\Complexesstack{\mathcal{C}\!\mathit{omplexes}}
% \Pic is the operator that assigns to X its picard group, usage \Pic(X)
% \Picardstack_{X/B} denotes the Picard stack of X over B
% \Picardfunctor_{X/B} denotes the Picard functor of X over B
\def\Pic{\mathop{\mathrm{Pic}}\nolimits}
\def\Picardstack{\mathcal{P}\!\mathit{ic}}
\def\Picardfunctor{\mathrm{Pic}}
\def\Deformationcategory{\mathcal{D}\!\mathit{ef}}

% --- appended 2026-10-12 (Chief): the project own definitions of the reference/slogan/history
% display environments, copied mechanically from the shared Stacks preamble lines 8-12
% (\usepackage{verbatim} plus the three \newenvironment lines, which the chapter preamble keeps
% verbatim from the project preamble). Nothing added, removed or redefined: these environments
% carry no mathematical content and the project itself typesets them as comments. Without these
% three lines a counted statement whose own text uses them fails with Environment slogan undefined.
\usepackage{verbatim}
\newenvironment{reference}{\comment}{\endcomment}
\newenvironment{slogan}{\comment}{\endcomment}
\newenvironment{history}{\comment}{\endcomment}
\usepackage{cleveref}
\crefname{originallemma}{original lemma}{original lemmas}
\Crefname{originallemma}{Original lemma}{Original lemmas}
\makeatletter
\newlabel{more-algebra-lemma-recognize-pseudo-coherent}{{064W}{source}{}{sourcepointer-1}{}}
\newlabel{more-algebra-lemma-recognize-pseudo-coherent@cref}{{[originallemma][0][]064W}{[1][0][]source}{}{}{}}
\makeatother
\crefname{originallemma}{original lemma}{original lemmas}
\Crefname{originallemma}{Original lemma}{Original lemmas}
\makeatletter
\newlabel{lemma-n-pseudo-module}{{064T}{source}{}{sourcepointer-2}{}}
\newlabel{lemma-n-pseudo-module@cref}{{[originallemma][0][]064T}{[1][0][]source}{}{}{}}
\makeatother
\crefname{originallemma}{original lemma}{original lemmas}
\Crefname{originallemma}{Original lemma}{Original lemmas}
\makeatletter
\newlabel{derived-remark-truncation-distinguished-triangle}{{08J5}{source}{}{sourcepointer-3}{}}
\newlabel{derived-remark-truncation-distinguished-triangle@cref}{{[originallemma][0][]08J5}{[1][0][]source}{}{}{}}
\makeatother
\crefname{originallemma}{original lemma}{original lemmas}
\Crefname{originallemma}{Original lemma}{Original lemmas}
\makeatletter
\newlabel{lemma-cone-pseudo-coherent}{{064R}{source}{}{sourcepointer-4}{}}
\newlabel{lemma-cone-pseudo-coherent@cref}{{[originallemma][0][]064R}{[1][0][]source}{}{}{}}
\makeatother
\title{Finiteness of the top cohomology modules recognises pseudo-coherence}
\author{The Stacks Project Authors}
\date{stacks-project:more-algebra.tex:local-snapshot-20261012; GNU FDL 1.2 or later}
\begin{document}
\maketitle
\setcounter{subsection}{1}
\noindent\textit{Editorial note.} All delivered mathematical text is the author's own, copied from the pinned Stacks Project chapter file more-algebra.tex; no statement or proof step was rewritten, reordered or removed.\par\medskip
\noindent\textit{Editorial note.} The chapter file pinned-original-source.tex carries only \textbackslash{}input\{preamble\} on its first line, so no preamble text lives inside the chapter. The retained file sources/stacks-more-algebra/retained-preamble.tex is a mechanical, line-for-line copy of the project's single shared preamble as staged in the snapshot (Xy-pic setup plus theorem declarations and mathematical macros). Nothing was added, removed or redefined, and the excerpt serves only as preamble, never as a span coordinate.\par\medskip
\noindent\textit{Editorial note.} Both counted spans start at the \textbackslash{}begin line of their environment, so the statement and the author's proof are delivered whole; surrounding prose of the section is deliberately outside the counted spans and is carried as standard prerequisites instead.\par\medskip
\noindent\textit{Editorial note.} The citation fields (source\_id, source\_version, authors, source\_url, license.name, source\_title) are identical to those of every other card this lane has produced from this chapter, so the batch selector treats them as one source.\par\medskip
\noindent\textit{Allowed standard prerequisites (editorial).}\begin{itemize}
\item Pseudo-coherence at a degree, tested through a finite free module mapping to the complex with the prescribed cohomology behaviour. defined earlier in the same section, source lines 16002-16156. Supplies the goal and the form of the witness constructed at the end.
\item A finite module placed in a single degree is pseudo-coherent in the relevant degree. the author cites Lemma lemma-n-pseudo-module at proof lines 16177-16179. Makes the high truncation pseudo-coherent.
\item Pseudo-coherence is stable under the operation taking a complex to the cone of a map, and the truncation triangle relates the two truncations to the complex. the author cites Lemma lemma-cone-pseudo-coherent and the truncation distinguished triangle at proof lines 16180-16186. Reduces the statement to the low truncation.
\end{itemize}\medskip
\ifdefined\hypertarget\hypertarget{sourcepointer-1}{}\fi
\noindent\textit{Source reference.} Original lemma 064W: focused-seeds/AGP/sources/stacks-snapshot-20261011/tags.txt, source lines 7880--7880. The lemma's own canonical tag, quoted so that readers can resolve the label inside the delivered text; it points at the counted statement itself and carries no proof content.\par\medskip
\ifdefined\hypertarget\hypertarget{sourcepointer-2}{}\fi
\noindent\textit{Source reference.} Original lemma 064T: focused-seeds/AGP/sources/stacks-snapshot-20261011/tags.txt, source lines 7877--7877. Navigation pointer to a statement the counted proof refers to but does not contain; nothing load-bearing is delegated to it.\par\medskip
\ifdefined\hypertarget\hypertarget{sourcepointer-3}{}\fi
\noindent\textit{Source reference.} Original lemma 08J5: focused-seeds/AGP/sources/stacks-snapshot-20261011/tags.txt, source lines 10838--10838. Navigation pointer to a statement the counted proof refers to but does not contain; nothing load-bearing is delegated to it.\par\medskip
\ifdefined\hypertarget\hypertarget{sourcepointer-4}{}\fi
\noindent\textit{Source reference.} Original lemma 064R: focused-seeds/AGP/sources/stacks-snapshot-20261011/tags.txt, source lines 7875--7875. Navigation pointer to a statement the counted proof refers to but does not contain; nothing load-bearing is delegated to it.\par\medskip
\begin{lemma}
\label{lemma-recognize-pseudo-coherent}
Let $R$ be a ring. Let $K^\bullet$ be a complex of $R$-modules.
Let $m \in \mathbf{Z}$.
\begin{enumerate}
\item If $H^i(K^\bullet) = 0$ for all $i \geq m$, then
$K^\bullet$ is $m$-pseudo-coherent.
\item If $H^i(K^\bullet) = 0$ for $i > m$ and $H^m(K^\bullet)$ is
a finite $R$-module, then $K^\bullet$ is $m$-pseudo-coherent.
\item If $H^i(K^\bullet) = 0$ for $i > m + 1$, the module
$H^{m + 1}(K^\bullet)$ is of finite presentation, and
$H^m(K^\bullet)$ is of finite type, then $K^\bullet$ is
$m$-pseudo-coherent.
\end{enumerate}
\end{lemma}

\begin{proof}
It suffices to prove (3). Set $M = H^{m + 1}(K^\bullet)$.
Note that $\tau_{\geq m + 1}K^\bullet$ is quasi-isomorphic to
$M[- m - 1]$. By
Lemma \ref{lemma-n-pseudo-module}
we see that $M[- m - 1]$ is $m$-pseudo-coherent. Since we have
the distinguished triangle
$$
(\tau_{\leq m}K^\bullet, K^\bullet, \tau_{\geq m + 1}K^\bullet)
$$
(Derived Categories, Remark
\ref{derived-remark-truncation-distinguished-triangle}) by
Lemma \ref{lemma-cone-pseudo-coherent}
it suffices to prove that $\tau_{\leq m}K^\bullet$ is pseudo-coherent.
By assumption $H^m(\tau_{\leq m}K^\bullet)$ is a finite type $R$-module.
Hence we can find a finite free $R$-module $E$ and a map
$E \to \Ker(d_K^m)$ such that the composition
$E \to \Ker(d_K^m) \to H^m(\tau_{\leq m}K^\bullet)$ is surjective.
Then $E[-m] \to \tau_{\leq m}K^\bullet$ witnesses the fact
that $\tau_{\leq m}K^\bullet$ is $m$-pseudo-coherent.
\end{proof}

\end{document}
